\documentclass[11pt]{article}
\usepackage[margin=1in]{geometry}
\usepackage{amsmath,amssymb,amsthm,mathtools}
\usepackage{microtype}
\usepackage{hyperref}
\usepackage{xcolor}
\usepackage{listings}
\usepackage{enumitem}
\usepackage{booktabs}

\hypersetup{
  colorlinks=true,
  linkcolor=blue!50!black,
  citecolor=blue!50!black,
  urlcolor=blue!50!black
}

\lstset{
  basicstyle=\ttfamily\scriptsize,
  breaklines=true,
  columns=fullflexible,
  frame=single,
  keepspaces=true,
  showstringspaces=false,
  tabsize=2,
  language=Python
}

\newtheorem{theorem}{Theorem}[section]
\newtheorem{lemma}[theorem]{Lemma}
\newtheorem{proposition}[theorem]{Proposition}
\newtheorem{corollary}[theorem]{Corollary}
\theoremstyle{definition}
\newtheorem{definition}[theorem]{Definition}
\newtheorem{remark}[theorem]{Remark}

\newcommand{\Hh}{\mathbb{H}}
\newcommand{\Rr}{\mathbb{R}}
\newcommand{\V}{V}
\newcommand{\Bn}{B_n}
\newcommand{\Sn}{S^n_0V}
\newcommand{\Ker}{\operatorname{ker}}
\newcommand{\rank}{\operatorname{rank}}
\newcommand{\Hom}{\operatorname{Hom}}
\newcommand{\ImH}{\operatorname{Im}\mathbb{H}}
\newcommand{\id}{\operatorname{id}}

\title{Residual Filtration and Probe-Depth Response\\
\large in Quaternionic Boundary Words\\[0.5em]
\normalsize A Companion Note to \emph{Free Numbers v0.2-pre.2}}
\author{Free Number Authors}
\date{2026-06-29}

\begin{document}
\maketitle

\begin{abstract}
This companion note studies residual kernels that arise from quaternionic boundary word spaces.
Let $V=\operatorname{Im}\mathbb{H}$ and $B_n=V^{\otimes n}$.
The reversed compression map
\[
  m_n(a_1|\cdots|a_n)=a_n\cdots a_1
\]
collapses boundary words into the quaternion algebra $\mathbb{H}$.
A single compression map cannot simultaneously serve as both a local closure judge and a residual-history witness.
To separate these roles, we introduce insertion-response profiles and probe-depth responses.

At length $2$, we prove
\[
  \Ker m_2=S^2_0V,
  \qquad
  K_2^{\mathrm{null}}=0,
\]
because the internal insertion response sends $S\in S^2_0V$ to $-2S$.
At length $3$, representation theory and an exact integer rank computation give
\[
  K_3^{\mathrm{null}}=S^3_0V.
\]
Thus the highest-spin cubic component is invisible at probe depth $1$.
We then define a depth-$2$ insertion profile and an internal-internal component $A_3$.
For $S\in S^3_0V$, exact computation gives
\[
  A_3(S)=4C_S,
\]
where $C_S(x,y)=S(x,y)$ is the ordinary two-variable contraction, viewed as a map into $\mathbb{H}$.
Hence the length-$3$ next-null residual is not absolutely invisible; it is invisible only at insufficient probe depth.
The general pattern suggested by these computations is not claimed as a theorem here.
\end{abstract}

\tableofcontents

\section{Introduction}

The main \emph{Free Numbers} paper introduces local quaternionic closure and global residual chart normalization.
Its central design principle is that local algebraic closure should not imply global collapse.
The present note does not replace that core construction.
Instead, it isolates a smaller technical mechanism that appears inside the residual side of the theory: boundary word compression, non-aligned kernels, and probe-depth response.

The guiding question is the following.
When a boundary word is killed by a compression map, has the residual history truly vanished, or has it only become invisible to that particular observation?

We show that the answer depends on probe depth.
At length $2$, the killed component is recovered by a depth-$1$ response.
At length $3$, the highest-spin component is invisible at depth $1$, but becomes visible at depth $2$.
This gives the first vertical step of a residual filtration indexed not only by word length $n$, but also by probe depth $d$.

The scope is intentionally narrow.
We prove only the length-$2$ and length-$3$ statements listed above and provide exact scripts for reproducing the finite computations.
We do not claim a general infinite ladder.

\paragraph{Relation to the core paper.}
The core paper constructs local closures $F_1\cong\mathbb{H}$ and a global non-collapsing chart structure.
Here we work in the associated pure-imaginary boundary word spaces
\[
  B_n=V^{\otimes n},\qquad V=\operatorname{Im}\mathbb{H},
\]
and study the finite linear maps produced by reversed quaternionic compression and insertion probes.

\section{Boundary words and reversed compression}

Let
\[
  V=\operatorname{Im}\mathbb{H}=\operatorname{span}_{\mathbb{R}}\{i,j,k\}.
\]
For pure imaginary quaternions $u,v\in V$, we use
\[
  uv=-u\cdot v+u\times v.
\]
Let
\[
  B_n=V^{\otimes n}
\]
be the real vector space spanned by boundary words
\[
  a_1|\cdots|a_n.
\]

\begin{definition}[Reversed compression]
The reversed compression map is
\[
  m_n:B_n\to\mathbb{H},
  \qquad
  m_n(a_1|\cdots|a_n)=a_n\cdots a_1.
\]
In particular,
\[
  m_2(u|v)=vu,
  \qquad
  m_3(a|b|c)=cba.
\]
\end{definition}

Let
\[
  K_n=\Ker m_n.
\]
The subspace $K_n$ records boundary combinations that vanish under the $n$-fold compression.
The central point of this note is that the vanishing of $m_n(x)$ does not necessarily mean that $x$ has no residual structure.
It means only that $x$ is invisible to the compression $m_n$.

\section{Why a single compression map is not enough}

A single compression map
\[
  m:B\to\mathbb{H}
\]
has two incompatible roles.
It can be used as a local closure judge, because it sends boundary words to quaternionic values.
But if one quotients by $\Ker m$, then all history lying inside $\Ker m$ has been collapsed.
Thus the same map cannot also serve as a witness for the residual history that it has killed.

This is the single-compression no-go principle:

\begin{quote}
A single compression map cannot simultaneously be the judge of local closure and the witness of residual persistence.
\end{quote}

To see residuals killed by $m_n$, one must use a non-aligned observation.
The insertion-response profiles below are the smallest such non-aligned witnesses.

\section{Depth-1 insertion response}

For $0\le r\le n$, define the insertion map
\[
  I_d^{(r)}(a_1|\cdots|a_n)
  =a_1|\cdots|a_r|d|a_{r+1}|\cdots|a_n.
\]
The $r$-th insertion response is
\[
  \mathcal{R}_n^{(r)}(x)(d)=m_{n+1}(I_d^{(r)}x).
\]
For a word $w=a_1|\cdots|a_n$, write
\[
  P_r(w)=a_r\cdots a_1,
  \qquad
  Q_r(w)=a_n\cdots a_{r+1}.
\]
Then
\[
  m_{n+1}(I_d^{(r)}w)=Q_r(w)dP_r(w).
\]

If $x\in\Ker m_n$, the exterior responses vanish:
\[
  \mathcal{R}_n^{(0)}(x)(d)=m_n(x)d=0,
  \qquad
  \mathcal{R}_n^{(n)}(x)(d)=dm_n(x)=0.
\]
Thus, at depth $1$, the internal insertions are the essential residual probes.

Define
\[
  K_n^{\mathrm{null}}
  =
  \Ker m_n\cap\Ker \mathcal{R}_n,
\]
where $\Ker \mathcal{R}_n$ means vanishing under all internal depth-$1$ insertion responses.

\section{Length 2: residual decomposition}

We identify $B_2=V\otimes V$.
As an $SO(3)$ representation,
\[
  V\otimes V=\mathbb{R}g\oplus\Lambda^2V\oplus S^2_0V.
\]
The dimensions are $9=1+3+5$.
Here $\mathbb{R}g$ is the trace component, $\Lambda^2V$ is the antisymmetric component, and $S^2_0V$ is the symmetric trace-free component.

\begin{lemma}
\label{lem:kerm2}
For $m_2(u|v)=vu$,
\[
  \Ker m_2=S^2_0V.
\]
\end{lemma}

\begin{proof}
Let
\[
  x=\sum_{a,b}T_{ab}e_a|e_b,
  \qquad
  e_1=i,
  \quad e_2=j,
  \quad e_3=k.
\]
Then
\[
  m_2(x)=\sum_{a,b}T_{ab}e_be_a.
\]
Using $e_be_a=-\delta_{ab}+e_b\times e_a$, we get
\[
  m_2(x)=-\operatorname{tr}(T)+\sum_{a,b}T_{ab}(e_b\times e_a).
\]
The real part sees the trace of $T$.
The vector part sees exactly the antisymmetric part of $T$, and the cross product identifies $\Lambda^2V$ with $V$.
Therefore $m_2(x)=0$ exactly when $T$ is symmetric and trace-free.
Hence $\Ker m_2=S^2_0V$.
\end{proof}

\begin{lemma}
\label{lem:r2}
For $S\in S^2_0V$,
\[
  \mathcal{R}_2^{(1)}(S)(c)=-2S(c).
\]
Thus
\[
  \mathcal{R}_2^{(1)}|_{S^2_0V}=-2\operatorname{id}.
\]
\end{lemma}

\begin{proof}
Write
\[
  S=\sum_{a,b}S_{ab}e_a|e_b,
  \qquad S_{ab}=S_{ba},
  \qquad \sum_a S_{aa}=0.
\]
The internal response is
\[
  \mathcal{R}_2^{(1)}(S)(c)=\sum_{a,b}S_{ab}e_bce_a.
\]
For pure imaginary $u,c,v\in V$, use the triple-product identity
\[
  ucv=-(u\cdot c)v-\det(u,c,v)+c(u\cdot v)-u(c\cdot v).
\]
Apply this with $u=e_b$ and $v=e_a$.
The determinant term vanishes after summing against the symmetric tensor $S_{ab}$.
The trace term vanishes because $S$ is trace-free.
The two remaining terms are both $-S(c)$.
Therefore
\[
  \mathcal{R}_2^{(1)}(S)(c)=-2S(c).
\]
\end{proof}

\begin{corollary}
\label{cor:k2null}
At length $2$,
\[
  K_2^{\mathrm{null}}=0.
\]
\end{corollary}

\section{Length 3: the first next-null residual}

At length $3$, the representation decomposition is
\[
  V^{\otimes3}\simeq V_0\oplus3V_1\oplus2V_2\oplus V_3.
\]
The dimensions are
\[
  1+3\cdot3+2\cdot5+7=27.
\]
The highest-spin component is
\[
  V_3=S^3_0V.
\]

The compression map
\[
  m_3:V^{\otimes3}\to\mathbb{H}
\]
has codomain $\mathbb{H}\simeq V_0\oplus V_1$.
The depth-$1$ response maps have codomain
\[
  \Hom(V,\mathbb{H})\simeq V_0\oplus2V_1\oplus V_2.
\]
There is no $V_3$ component in either codomain.
By equivariance and Schur's lemma,
\[
  S^3_0V\subseteq K_3^{\mathrm{null}}.
\]

To prove equality, define
\[
  T_3=m_3\oplus\mathcal{R}_3^{(1)}\oplus\mathcal{R}_3^{(2)}.
\]
This is a map
\[
  T_3:V^{\otimes3}\to
  \mathbb{H}\oplus\Hom(V,\mathbb{H})\oplus\Hom(V,\mathbb{H}).
\]
The domain has dimension $27$ and the codomain has dimension $4+12+12=28$.
In the basis $e_a|e_b|e_c$, the matrix entries are determined by
\[
  m_3(e_a|e_b|e_c)=e_ce_be_a,
\]
\[
  \mathcal{R}_3^{(1)}(e_a|e_b|e_c)(d)=e_ce_bde_a,
\]
and
\[
  \mathcal{R}_3^{(2)}(e_a|e_b|e_c)(d)=e_cde_be_a.
\]
Exact integer computation gives
\[
  \rank(T_3)=20.
\]
Thus
\[
  \dim\Ker T_3=27-20=7.
\]
Since $\dim S^3_0V=7$, we get the equality.

\begin{theorem}[Length-3 next-null residual]
\label{thm:k3null}
At length $3$,
\[
  K_3^{\mathrm{null}}=S^3_0V.
\]
\end{theorem}

This means that the highest-spin cubic component is invisible to the depth-$1$ response.
The next section shows that it is nevertheless visible at probe depth $2$.

\section{Probe depth 2: the first vertical step}

The full depth-$2$ insertion profile consists of all maps obtained by inserting two probe variables into a length-$3$ boundary word and then applying the length-$5$ reversed compression $m_5$.

We use one component of that full profile.
For a word $e_a|e_b|e_c$, take the internal-internal double insertion
\[
  e_a|x|e_b|y|e_c.
\]
After reversed compression, this gives
\[
  m_5(e_a|x|e_b|y|e_c)=e_cye_bxe_a.
\]
Define
\[
  A_3(T)(x,y)=\sum_{a,b,c}T_{abc}e_cye_bxe_a
\]
for
\[
  T=\sum_{a,b,c}T_{abc}e_a|e_b|e_c.
\]

\begin{remark}[Component versus full profile]
The map $A_3$ is not asserted to be the whole depth-$2$ profile.
It is one internal-internal component of the canonical full profile.
We do not claim that other depth-$2$ insertion components vanish.
We use only the safe implication that if one component is injective on a subspace, then the full profile detects that subspace.
\end{remark}

For $S\in S^3_0V$, define the contraction embedding
\[
  C_S:V\otimes V\to V\subset\mathbb{H},
  \qquad
  C_S(x,y)=S(x,y).
\]
In coordinates,
\[
  C_S(x,y)=\sum_{a,b,c}S_{abc}x_ay_be_c.
\]
Thus
\[
  C:S^3_0V\to\Hom(V^{\otimes2},\mathbb{H})
\]
is the natural contraction embedding.

\subsection{Why depth 2 can see spin 3}

At depth $1$,
\[
  \Hom(V,\mathbb{H})\simeq V_0\oplus2V_1\oplus V_2,
\]
so there is no $V_3$ component.
At depth $2$,
\[
  \Hom(V^{\otimes2},\mathbb{H})
  \simeq
  2V_0\oplus4V_1\oplus3V_2\oplus V_3.
\]
The depth-$2$ codomain therefore contains $V_3$ for the first time.
This does not by itself prove visibility, but it makes visibility representation-theoretically possible.
The coefficient is computed directly.

\begin{theorem}[Depth-2 detection of the cubic highest-spin residual]
\label{thm:a3}
For every $S\in S^3_0V$,
\[
  A_3(S)=4C_S.
\]
Equivalently, for all $x,y\in V$,
\[
  A_3(S)(x,y)=4S(x,y).
\]
Consequently, $A_3$ is injective on $S^3_0V$.
Since $A_3$ is one component of the full depth-$2$ insertion profile, the full depth-$2$ profile detects the entire space $K_3^{\mathrm{null}}=S^3_0V$.
\end{theorem}

\begin{proof}
The identity is verified by exact integer computation on a basis of $S^3_0V$.
The script in Appendix~\ref{app:a3} constructs the map $A_3$, constructs the contraction map $C$, constructs the symmetric trace-free subspace $S^3_0V$, and checks that $(A_3-4C)$ vanishes on a basis.
It also checks that the restricted image has rank $7$.
\end{proof}

\subsection{Example}

Consider
\[
  S=i|i|i-i|j|j-j|i|j-j|j|i.
\]
This is the harmonic cubic component corresponding to $x^3-3xy^2$, hence $S\in S^3_0V$.
From Theorem~\ref{thm:k3null}, it is invisible at depth $1$.
At depth $2$,
\[
  A_3(S)(i,i)=4i,
\]
\[
  A_3(S)(i,j)=-4j,
\]
and
\[
  A_3(S)(j,j)=-4i.
\]
These agree with $4S(i,i)$, $4S(i,j)$, and $4S(j,j)$.

\section{Consequences}

Theorem~\ref{thm:k3null} says that $S^3_0V$ is invisible at depth $1$.
Theorem~\ref{thm:a3} says that the same residual is visible at depth $2$.
Thus the residual is not absolutely invisible.
It is invisible only at insufficient probe depth.

This is the first vertical step of the residual filtration:
\[
  \text{depth }1:\ \text{invisible},
  \qquad
  \text{depth }2:\ \text{visible}.
\]
Together with the length-$2$ result, the picture is:
\[
  K_2^{\mathrm{null}}=0,
  \qquad
  K_3^{\mathrm{null}}=S^3_0V,
  \qquad
  A_3|_{S^3_0V}=4C.
\]

\section{Non-claims and suggested pattern}

The computations above suggest the pattern
\[
  A_n|_{S^n_0V}=(-2)^{n-1}\operatorname{id},
\]
with coefficient $-2$ at $n=2$ and $+4=(-2)^2$ at $n=3$.
However, this paper does not claim the general theorem.
It does not claim that all highest-spin residuals are detected exactly at depth $n-1$.
It does not claim an infinite residual ladder.

The only claims established here are:
\[
  K_2^{\mathrm{null}}=0,
  \qquad
  K_3^{\mathrm{null}}=S^3_0V,
  \qquad
  A_3(S)=4C_S\quad(S\in S^3_0V).
\]

The general spin-depth pattern remains a suggested direction for later work.

\section{Conclusion}

Residual kernels in quaternionic boundary word spaces are not merely artifacts of failed compression.
They carry information that may become visible at higher probe depth.
At length $2$, depth $1$ already recovers the killed symmetric trace-free component.
At length $3$, depth $1$ fails to see the highest-spin cubic component, but depth $2$ detects it as $4C_S$.

This establishes a small but precise two-axis picture: word length controls which highest-spin residuals appear, while probe depth controls when they become visible.
The resulting structure is a residual filtration in both length and probe depth.

\appendix

\section{Exact rank check for the length-3 null residual}
\label{app:t3}

The following script constructs the exact integer matrix $T_3$ and computes its rank.

\begin{lstlisting}
import sympy as sp

# Quaternion basis order:
# 0 = 1, 1 = i, 2 = j, 3 = k

def qmul(x, y):
    a0, a1, a2, a3 = x
    b0, b1, b2, b3 = y
    return [
        a0*b0 - a1*b1 - a2*b2 - a3*b3,
        a0*b1 + a1*b0 + a2*b3 - a3*b2,
        a0*b2 - a1*b3 + a2*b0 + a3*b1,
        a0*b3 + a1*b2 - a2*b1 + a3*b0,
    ]

basis = [
    [1, 0, 0, 0],  # 1
    [0, 1, 0, 0],  # i
    [0, 0, 1, 0],  # j
    [0, 0, 0, 1],  # k
]

imag = [basis[1], basis[2], basis[3]]

def prod(seq):
    v = basis[0]
    for x in seq:
        v = qmul(v, x)
    return v

cols = []

for a in range(3):
    for b in range(3):
        for c in range(3):
            ea, eb, ec = imag[a], imag[b], imag[c]
            col = []

            # m_3(e_a|e_b|e_c) = e_c e_b e_a
            col += prod([ec, eb, ea])

            # R_3^(1)(e_a|e_b|e_c)(d) = e_c e_b d e_a
            for d in range(3):
                ed = imag[d]
                col += prod([ec, eb, ed, ea])

            # R_3^(2)(e_a|e_b|e_c)(d) = e_c d e_b e_a
            for d in range(3):
                ed = imag[d]
                col += prod([ec, ed, eb, ea])

            cols.append(col)

M = sp.Matrix.hstack(*[sp.Matrix(col) for col in cols])

print(M.shape)
print(M.rank())
\end{lstlisting}

Expected output:
\begin{lstlisting}
(28, 27)
20
\end{lstlisting}

\section{Exact check for the depth-2 identity}
\label{app:a3}

The following script verifies $A_3(S)=4C_S$ on a basis of $S^3_0V$.

\begin{lstlisting}
import sympy as sp
from itertools import product, permutations

# Quaternion basis order:
# 0 = 1, 1 = i, 2 = j, 3 = k

def qmul(x, y):
    a0, a1, a2, a3 = x
    b0, b1, b2, b3 = y
    return sp.Matrix([
        a0*b0 - a1*b1 - a2*b2 - a3*b3,
        a0*b1 + a1*b0 + a2*b3 - a3*b2,
        a0*b2 - a1*b3 + a2*b0 + a3*b1,
        a0*b3 + a1*b2 - a2*b1 + a3*b0,
    ])

basis = [
    sp.Matrix([1, 0, 0, 0]),  # 1
    sp.Matrix([0, 1, 0, 0]),  # i
    sp.Matrix([0, 0, 1, 0]),  # j
    sp.Matrix([0, 0, 0, 1]),  # k
]

imag = [basis[1], basis[2], basis[3]]

def prodq(seq):
    v = basis[0]
    for x in seq:
        v = qmul(v, x)
    return v

triples = list(product(range(3), repeat=3))
idx = {t: i for i, t in enumerate(triples)}

# A3 matrix:
# rows are indexed by (x, y, quaternion_component), dimension 3*3*4 = 36.
# columns are indexed by (a, b, c), dimension 27.
A = sp.zeros(36, 27)

# C matrix:
# ordinary contraction C_S(x,y)=S(x,y), viewed in H.
C = sp.zeros(36, 27)

for col, (a, b, c) in enumerate(triples):
    ea, eb, ec = imag[a], imag[b], imag[c]

    for x in range(3):
        ex = imag[x]
        for y in range(3):
            ey = imag[y]
            rowbase = (x*3 + y)*4

            # A3(e_a|e_b|e_c)(x,y) = e_c y e_b x e_a
            out = prodq([ec, ey, eb, ex, ea])
            for h in range(4):
                A[rowbase + h, col] = out[h]

            # C(e_a|e_b|e_c)(x,y) = delta_{a,x} delta_{b,y} e_c
            if a == x and b == y:
                for h in range(4):
                    C[rowbase + h, col] = ec[h]

# Build equations for S^3_0 V inside V^{tensor 3}.
eqs = []

# Symmetry equations: T_abc = T_perm(a,b,c).
for t in triples:
    for p in set(permutations(t)):
        if t < p:
            row = [0]*27
            row[idx[t]] = 1
            row[idx[p]] = -1
            eqs.append(row)

# Trace-free equations: sum_a T_aac = 0 for c = 0,1,2.
for c in range(3):
    row = [0]*27
    for a in range(3):
        row[idx[(a, a, c)]] = 1
    eqs.append(row)

E = sp.Matrix(eqs)
S3_basis = E.nullspace()

print("dim S^3_0 V =", len(S3_basis))

D = A - 4*C

checks = [D*v == sp.zeros(36, 1) for v in S3_basis]
print("A3 - 4C vanishes on basis:", all(checks))

restricted_images = sp.Matrix.hstack(*[A*v for v in S3_basis])
print("rank of A3 on S^3_0 V =", restricted_images.rank())
\end{lstlisting}

Expected output:
\begin{lstlisting}
dim S^3_0 V = 7
A3 - 4C vanishes on basis: True
rank of A3 on S^3_0 V = 7
\end{lstlisting}

\end{document}
